% Lokalni simboli predavanja 01, usklađeni sa stilom B.
% Dvoulazni simbol: vrsta (and/or/xor/nand/nor/xnor), ulazi (nn/ii), izlaz.
\newcommand{\DEkapija}[3]{%
  \begin{tikzpicture}[circuit logic US,DEoznaka,baseline=(g.center)]
    \node[#1 gate US,draw,logic gate inputs=#2,minimum width=12mm,minimum height=10mm] (g) {};
    \draw[DEvod] (g.input 1) -- ++(-.35,0) node[left] {$x$};
    \draw[DEvod] (g.input 2) -- ++(-.35,0) node[left] {$y$};
    \draw[DEvod] (g.output) -- ++(.35,0) node[right] {$#3$};
  \end{tikzpicture}}
\newcommand{\DEtroXor}[1]{%
  \begin{tikzpicture}[DEoznaka,baseline=(current bounding box.center)]
    \draw[DEvod] (-.55,.6) .. controls (.05,.6) and (.55,.4) .. (.85,0)
      .. controls (.55,-.4) and (.05,-.6) .. (-.55,-.6)
      .. controls (-.12,-.2) and (-.12,.2) .. cycle;
    \draw[DEvod] (-.7,.6) .. controls (-.27,.2) and (-.27,-.2) .. (-.7,-.6);
    \draw[DEvod] (-1.1,.36) node[left] {$x_1$} -- (-.37,.36);
    \draw[DEvod] (-1.1,0) node[left] {$x_2$} -- (-.23,0);
    \draw[DEvod] (-1.1,-.36) node[left] {$x_3$} -- (-.37,-.36);
    \ifnum#1=1
      \draw[DEvod] (.92,0) circle (.07);
      \draw[DEvod] (.99,0)--(1.4,0) node[right] {$y$};
    \else
      \draw[DEvod] (.85,0)--(1.4,0) node[right] {$y$};
    \fi
  \end{tikzpicture}}
\newcommand{\DEtroulaz}[1]{%
  \ifstrequal{#1}{xor}{\DEtroXor{0}}{\ifstrequal{#1}{xnor}{\DEtroXor{1}}{%
  \begin{tikzpicture}[circuit logic US,DEoznaka,baseline=(g.center)]
    \node[#1 gate US,draw,logic gate inputs=nnn,minimum width=12mm,minimum height=12mm] (g) {};
    \foreach \i in {1,2,3}
      \draw[DEvod] (g.input \i) -- ++(-.35,0) node[left] {$x_\i$};
    \draw[DEvod] (g.output) -- ++(.35,0) node[right] {$y$};
  \end{tikzpicture}}}}
% Tri ulaza kroz dva dvoulazna kola; opcija 1 dodaje izlazni invertor.
\newcommand{\DEkaskada}[3][0]{%
  \begin{tikzpicture}[circuit logic US,DEoznaka,baseline=(b.center)]
    \node[#2 gate US,draw,logic gate inputs=nn,minimum width=11mm,minimum height=10mm] (a) at (0,.35) {};
    \node[#3 gate US,draw,logic gate inputs=nn,minimum width=11mm,minimum height=10mm] (b) at (1.7,-.3) {};
    \draw[DEvod] (a.input 1)--++(-.4,0) node[left] {$x_1$};
    \draw[DEvod] (a.input 2)--++(-.4,0) node[left] {$x_2$};
    \draw[DEvod] (a.output)--++(.2,0)|-(b.input 1);
    \draw[DEvod] (b.input 2)--++(-.4,0) node[left] {$x_3$};
    \ifnum#1=1
      \node[not gate US,draw,minimum width=8mm] (c) at (3.2,-.3) {};
      \draw[DEvod] (b.output)--(c.input);
      \draw[DEvod] (c.output)--++(.35,0) node[right] {$y$};
    \else
      \draw[DEvod] (b.output)--++(.35,0) node[right] {$y$};
    \fi
  \end{tikzpicture}}
% Kaskada paritetnih kola: ulaz x_3 i međusignal imaju odvojene trase.
% Dva tipa omogućavaju i ekvivalentnu realizaciju troulaznog EXNILI kola.
\newcommand{\DEparitetKaskadaDva}[2]{%
  \begin{tikzpicture}[circuit logic US,DEoznaka,baseline=(current bounding box.center)]
    \node[#1 gate US,draw,logic gate inputs=nn,minimum width=11mm,minimum height=10mm] (prvo) at (0,.65) {};
    \node[#2 gate US,draw,logic gate inputs=nn,minimum width=11mm,minimum height=10mm] (drugo) at (2.55,-.45) {};
    \draw[DEvod] (prvo.input 1) -- ++(-.4,0) node[left] {$x_1$};
    \draw[DEvod] (prvo.input 2) -- ++(-.4,0) node[left] {$x_2$};
    \coordinate (spust) at ($(prvo.output)+(.35,0)$);
    \draw[DEvod] (prvo.output) -- (spust) |- (drugo.input 1);
    \draw[DEvod] (drugo.input 2) -- ++(-.5,0) node[left] {$x_3$};
    \draw[DEvod] (drugo.output) -- ++(.35,0) node[right] {$y$};
  \end{tikzpicture}}
\newcommand{\DEparitetKaskada}[1]{\DEparitetKaskadaDva{#1}{#1}}
